\subsection{The Saad \& Solla Equations}
Compressing the problem to order parameters $(Q, M, \vec{a})$ is only half the
story.The other half is the \emph{high-dimensional limit} $d \to \infty$: in this regime,
the stochastic fluctuations of the order parameters around their mean become
negligible, and the entire training trajectory concentrates on a deterministic path
governed by a system of \emph{ordinary differential equations} (ODEs).

\paragraph{Low dimensional dynamics for an high-dimensional problem.}
We can make Equation~\eqref{eq:intro:sgd} explicit in terms of the local fields, and write the stochastic dynamics of the weights $\vec{w}_i$ and the output weights $a_i$ of the student network as
\begin{align}
    \vec{w}_j^{\nu+1} &= \vec{w}_j^{\nu} + \frac{\gamma}{p\sqrt{d}} \left( y(\vec{\lambda}^{\star,\nu}) - h(\vec{\lambda}^{\nu}) \right) a_j \sigma'(\lambda^\nu_j) \vec{x}^\nu, \\
    %
    a_j^{\nu+1} &= a_j^{\nu} + \frac{\gamma}{pd} \left( y(\vec{\lambda}^{\star,\nu}) - h(\vec{\lambda}^{\nu}) \right) \sigma(\lambda^\nu_j)
\end{align}
where \(\vec{\lambda}^{\nu} = \frac{1}{\sqrt{d}}(\langle \vec{w}_1, \vec{x}^\nu \rangle, \ldots, \langle \vec{w}_p, \vec{x}^\nu \rangle)\) and \(\vec{\lambda}^{\star,\nu} = \frac{1}{\sqrt{d}}(\langle \vec{w}_1^\star, \vec{x}^\nu \rangle, \ldots, \langle \vec{w}_k^\star, \vec{x}^\nu \rangle)\) are the \emph{local fields} (i.e., pre-activations) of the student and teacher networks, respectively, relative to the current input $\vec{x}^\nu$ used in the $\nu$-th SGD step. The quantity \(y^\nu(\lambda^{\star,\nu}) \coloneqq y(\vec{\lambda}^{\star,\nu}) + \sqrt{\Delta}\xi^\nu\) is the noisy label of the current training sample.
The second-layer weights $a_j$ are updated with a rescaled learning rate $\gamma/d$, which keeps both layers evolving on the same timescale as $d \to \infty$.
With a bit of extra work we can derive the update rule for the order parameters
\begin{equation} \label{eq:intro:op-update}
\begin{aligned}
    Q^{\nu+1}_{jl} =& Q^{\nu}_{jl} + \frac{\gamma}{pd} \left( y(\vec{\lambda}^{\star,\nu}) - h(\vec{\lambda}^{\nu}) \right) \left( a_l \sigma'(\lambda^\nu_l) \lambda^\nu_j + a_j \sigma'(\lambda^\nu_j) \lambda^\nu_l  \right) + \\ & + \frac{\gamma}{p d} \frac{\gamma}{p}\left( y(\vec{\lambda}^{\star,\nu}) - h(\vec{\lambda}^{\nu}) \right)^2 a_j a_l \sigma'(\lambda^\nu_j) \sigma'(\lambda^\nu_l) \frac{\norm{\vec{x}^\nu}^2}{d}, \\
    M^{\nu+1}_{jr} =& M^{\nu}_{jr} + \frac{\gamma}{pd} \left( y(\vec{\lambda}^{\star,\nu}) - h(\vec{\lambda}^{\nu}) \right) a_j \sigma'(\lambda^\nu_j) \lambda^{\star,\nu}_r , \\
    a_j^{\nu+1} =& a_j^{\nu} + \frac{\gamma}{pd} \left( y(\vec{\lambda}^{\star,\nu}) - h(\vec{\lambda}^{\nu}) \right) \sigma(\lambda^\nu_j).
\end{aligned}
\end{equation}
Since we are interested in the high-dimensional limit, the factor $\sfrac{\norm{\vec{x}^\nu}^2}{d}$ in the $Q$ update can be replaced by its expectation $\mathbb{E}\left[\sfrac{\norm{\vec{x}^\nu}^2}{d}\right] = 1$. \addnote{This is the classical \emph{à la physique} argument from \citet{saad1995exact}. In the next section we will report a formalization of this argument.}
The update rules above are still stochastic, but they do not explicitly depend on the input $\vec{x}^\nu$ anymore, only on the local fields $\vec{\lambda}^\nu$ and $\vec{\lambda}^{\star,\nu}$, which are Gaussian random variables with covariance given by the order parameters themselves. In other words, we have a \emph{duality} between the process in the space of the weights \(W\in\mathbb{R}^{p\times d}\) and the process in the space of the order parameters \((Q, M, \vec{a})\in\mathbb{R}^{p\times p}\times\mathbb{R}^{p\times k}\times\mathbb{R}^p\)
\begin{align*}
    \text{update rule for } W \quad &\leftrightarrow \quad \text{update rule for } (Q, M, \vec{a}) \\
    \text{sample } \vec{x}^\nu \sim \mathcal{N}(0, I_d) \quad &\leftrightarrow \quad \text{sample } (\vec{\lambda}^\nu, \vec{\lambda}^{\star,\nu}) \sim \mathcal{N}(0, \Omega) \\
    \text{population risk } \mathcal{R}(W) \quad &\leftrightarrow \quad \text{population risk } \mathcal{R}(Q, M, \vec{a}).
\end{align*}
The duality allows us to have a full characterization of the random SGD process in dimension \(d\) that is independent of the \(d\) itself. The population risk is a quanity that fully characterizes the generalization properties of the student network, and it can be computed from the order parameters; on the other hand, the dynamics is self-contained in the order parameters space, so we can study the generalization properties of the student network without having to deal with the high-dimensional weights \(W\) directly.

\paragraph{The deterministic equations}
The update rules in Equation~\eqref{eq:intro:op-update} can be rearranged as
\begin{equation}
\begin{aligned}
    \frac{Q^{\nu+1}_{jl}-Q^{\nu}_{jl}}{\sfrac{\gamma}{pd}} =& \left( y(\vec{\lambda}^{\star,\nu}) - h(\vec{\lambda}^{\nu}) \right) \left( a_l \sigma'(\lambda^\nu_l) \lambda^\nu_j + a_j \sigma'(\lambda^\nu_j) \lambda^\nu_l  \right) + \\ & + \frac{\gamma}{p}\left( y(\vec{\lambda}^{\star,\nu}) - h(\vec{\lambda}^{\nu}) \right)^2 a_j a_l \sigma'(\lambda^\nu_j) \sigma'(\lambda^\nu_l), \\
    \frac{M^{\nu+1}_{jr}-M^{\nu}_{jr}}{\sfrac{\gamma}{pd}} =& \left( y(\vec{\lambda}^{\star,\nu}) - h(\vec{\lambda}^{\nu}) \right) a_j \sigma'(\lambda^\nu_j), \\
    \frac{a_j^{\nu+1}-a_j^{\nu}}{\sfrac{\gamma}{pd}} =& \left( y(\vec{\lambda}^{\star,\nu}) - h(\vec{\lambda}^{\nu}) \right) \sigma(\lambda^\nu_j).
\end{aligned}
\end{equation}
Intuitively, when \(d\to\infty\) the left-hand side of the equations above can be interpreted as a derivative with respect to a continuous time variable \(t = \nu \frac{\gamma}{pd}\), and the right hand side can be replaced by its expectation over the Gaussian distribution of the local fields, which is parametrized by the order parameters themselves. \addnote{In the traditional derivation of Saad\&Solla equations, the usual time sacling is $t=\frac{\nu}{d}$. Later will clarify why we prefer this scaling.} This leads to a system of ODEs for the order parameters, known as the \emph{Saad \& Solla equations}\cite{saad1995exact,saad1996dynamics}:
\begin{equation} \label{eq:intro:ss-odes} \tag{SS-ODEs}
\begin{aligned}
    \dod{Q_{jl}}{t} =& \overbrace{\mathbb{E}_{(\vec{\lambda}, \vec{\lambda}^\star) \sim \mathcal{N}(0, \Omega)}
    \left[ \left( h^\star(\vec{\lambda}^{\star}) - h(\vec{\lambda}) \right) \left( a_l \sigma'(\lambda_l) \lambda_j + a_j \sigma'(\lambda_j) \lambda_l  \right)\right]}^{\textstyle \eqqcolon \Phi^{Q}_{jl}(\Omega, a)}\\ & + \frac{\gamma}{p} \underbrace{\mathbb{E}_{(\vec{\lambda}, \vec{\lambda}^\star) \sim \mathcal{N}(0, \Omega)}\left[ \left(\left( h^\star(\vec{\lambda}^{\star}) - h(\vec{\lambda}) \right)^2+\Delta\right) a_j a_l \sigma'(\lambda_j) \sigma'(\lambda_l) \right]}_{\textstyle \eqqcolon \Psi^{Q}_{jl}(\Omega, a)}, \\[0.9em]
    \dod{M_{jr}}{t} =& \underbrace{\mathbb{E}_{(\vec{\lambda}, \vec{\lambda}^\star) \sim \mathcal{N}(0, \Omega)} \left[ \left( h^\star(\vec{\lambda}^{\star}) - h(\vec{\lambda}) \right) a_j \sigma'(\lambda_j) \lambda^{\star}_r \right]}_{\textstyle \eqqcolon \Phi^{M}_{jr}(\Omega, a)}, \\[0.9em]
    \dod{a_j}{t} =& \underbrace{\mathbb{E}_{(\vec{\lambda}, \vec{\lambda}^\star) \sim \mathcal{N}(0, \Omega)} \left[ \left( h^\star(\vec{\lambda}^{\star}) - h(\vec{\lambda}) \right) \sigma(\lambda_j) \right]}_{\textstyle \eqqcolon \Phi^{a}_{j}(\Omega, a)}.
\end{aligned}
\end{equation}
The trajectory drawn by these equations is completely determined by the initial conditions. After that the dynamics is deterministic, and the randomness of the training process is averaged out in the high-dimensional limit.
\begin{remark}
    Note that even though the Saad \& Solla determine a deterministic trajectory for the order parameters, the actual trajectory of the weights $W$ is still stochastic, and impossible to characterize without knowing the specific realization of the training data \(x^\nu\). The Saad \& Solla equations only describe the trajectory of the order parameters, and consequently the population risk.
\end{remark}

The formalization of the argument above, and the proof of the convergence of the stochastic process to the deterministic one was first given by \citet{goldt2019dynamics}. \addnote{ The original theorem was stated for matching student and teacher networks, but the proof can be easily extended to the more general case of a multi-index teacher, and network student.}
\begin{theorem}[Convergence of S\&S Equations, \cite{goldt2019dynamics}] \label{thm:intro:ss-convergence}
    Let's suppose that
  \begin{enumerate}
    \item The function \(\sigma\) is bounded and its derivatives up to the second order exist and are bounded, too;
    \item The function \(h^\star\) grows at most like a polynomial, i.e. \(|h^\star(\vec{\lambda}^\star)| \le K(1 + \|\vec{\lambda}^\star\|^p)\) for some constants \(K\) and \(p\);
    \item \(W^\star\) is of rank \(k\);
    \item The initial macroscopic state \({\Omega^0}\) is deterministic and bounded by a constant;
    \item The constants \(\Delta, p, k\) and \(\gamma\) are all finite.
  \end{enumerate}
  Let \(T>0\), then the overlap matrix satisfies 
  \begin{equation}
    \max_{0\le\nu\le Td} \E\left[
      \left\|({\Omega}^\nu,a^\nu) - \left(\tilde{{\Omega}}\left(\frac{\nu}{d}\right), \tilde{a}\left(\frac{\nu}{d}\right)\right)\right\|
    \right]
    \le
    \frac{C{(T)}}{\sqrt{d}},
  \end{equation}
  where \(\left(\tilde{{\Omega}}, \tilde{a}\right)\) it's the unique solution of the Equations~\eqref{eq:intro:ss-odes},
  with the initial condition \[\left(\tilde{{\Omega}}{(0)}, \tilde{a{(0)}}\right) = \left(\Omega^0, a^0\right).\]
\end{theorem}

\paragraph{The SGD noise.} It is useful trying to interpret the actual content of the Saad \& Solla equations; for this with follow the nomenclature use by \citet{veiga2022phase}. We have two typology of terms:
\begin{enumerate}
    \item \emph{learning terms} \(\Phi\): they are the terms leading the dynamic to align the student with the teacher, and consequently to reduce the population risk; it can be shown that thse same terms are the one remaining in the \emph{gradient flow} limit, thus driving the learning.
    \item \emph{variance term} \(\Psi\): they are the terms that arise from the stochasticity of the training process, and they are responsible for the fluctuations of the order parameters around their mean trajectory; it is interesting to note that they are proportional to the learning rate \(\gamma\), and the \emph{do not} vanish for \(\Delta = 0\). This term is the crucial noise introduced by SGD, fundamentally differentiating it from the gradient flow limit.
\end{enumerate}
The \emph{variance term} is responsible for the \emph{plateau} phenomenon \cite{veiga2022phase}, where the order parameters get stuck in a suboptimal configuration and can't reach the exact global minimum even in the limit \(t\to\infty\).
% TODO put a pictureof the plateau phenomenon.

Practitioners often address this problem by using a decaying learning rate, which effectively makes the variance term vanish, combined to larger batch sizes, which reduce the variance of the stochastic gradients by averaging the noise of the labels. 
It is natural to ask how the Equations~\eqref{eq:intro:ss-odes} change when we consider a learning rate and a batch size that are also scaling with the dimension \(d\); this is the exact topic of Chapter~\ref{chap:paper-2406}.
% add a couple of lines explaining the new saad and solla equation we found: intra-batch non-asymptotic corrections, phase diagram where we can kill variance, but learn with higher sample complexity, or be on the edge of nearning/not learning, and have the plateu. paper-2406

\paragraph{Examples}
The Equation~\eqref{eq:intro:ss-odes} can fully describe the dynamics only if it is actually possible to analyticaly integrate the expectations on the right-hand side, arriving to a closed-form expression for \(\Phi, \Psi\) as a function of the order parameters. This stronly depends on the choice of the activation function \(\sigma\) and the teacher function \(h^\star\); in literature the problem is often tackled by considering simple setting of \emph{matching architectures}.

Let's assume for simplicity that the \emph{teacher} is a two-layer network with $k$ hidden neurons and the same activation function $\sigma$, as the \emph{student} network is with $p$ hidden neurons
\begin{equation}
    f^\star(\vec{x}) = \frac1k \sum_{j=1}^k a_j^\star \sigma\left(\frac{\vec{w}_j^\star \cdot \vec{x}}{\sqrt{d}}\right)
    \quad\text{or equivalently}\quad
    h^\star(\vec{\lambda}^\star) = \frac1k \sum_{j=1}^k a_j^\star \sigma(\lambda_j^\star)
\end{equation}
The whole integrability of Equation~\eqref{eq:intro:ss-odes} boils down to the possibility to compute the following expectations given any generic covariance \(C\)
\begin{align}
    I_2^\text{risk}(C_2) &\coloneqq \mathbb{E}_{(z_1, z_2) \sim \mathcal{N}(0, C_2)}\left[ \sigma(z_1) \sigma(z_2) \right], \\
    I_2^\text{noise}(C_2) &\coloneqq \mathbb{E}_{(z_1, z_2) \sim \mathcal{N}(0, C_2)}\left[ \sigma'(z_1) \sigma'(z_2) \right], \\
    I_3(C_3) &\coloneqq \mathbb{E}_{(z_1, z_2, z_3) \sim \mathcal{N}(0, C_3)}\left[ \sigma(z_1) \sigma'(z_2) z_3 \right], \\
    I_4(C_4) &\coloneqq \mathbb{E}_{(z_1, z_2, z_3, z_4) \sim \mathcal{N}(0, C_4)}\left[ \sigma'(z_1) \sigma'(z_2) z_3 z_4 \right].
\end{align}
The term \(I_2^\text{risk}\) is the one that appears in the expression of the population risk; the term \(I_3\) is the one that appears in the learning terms \(\Phi\); the terms \(I_2^\text{noise}\) and \(I_4\) are the ones that appear in the variance term \(\Psi\).

In this case, a lot of literature has been devoted into discussing the possibility to have exact asymptotic equation for the dynamics depending on the choice of the activation function \(\sigma\) \addnote{Some of this activations do not fully satisfy the conditions of Theorem~\ref{thm:intro:ss-convergence}, but they have been proven to work empirically. On the other hand, the theorem is \emph{probably} not tight.}
\begin{itemize}
    \item \emph{linear} \(\sigma(x) = x\): this case is trivial, since the student and the teacher are just linear models, and the dynamics can be fully solved in closed form;
    \item \emph{error function} \(\sigma(x) = \erf(\sfrac{x}{\sqrt{2}})\): this is the original setting of \citet{saad1995exact}, where the equations can be fully solved in closed form, and the dynamics can be fully characterized;
    \item \emph{ReLU} \(\sigma(x) = \max(0, x)\): this is the most common activation function used in practice, but it does not have a closed form solution for the expectations \(I_4\). The other integrals have been studied in \citet{yoshida2019data} using result from \citet{yoshida2017statistical};
    \item \emph{phase retrieval} \(\sigma(x) = x^2\): the study of Equations~\eqref{eq:intro:ss-odes} for this activation function is the full content of Chapter~\ref{chap:paper-2305}. We determined the exact closed form of the expectations, and used the result to study the number of samples needed to learn the teacher network.
    \item \emph{polynomial} \(\sigma(x) = p(x)\): the framework introduced in Chapter~\ref{chap:paper-2305} can be extended to any polynomial activation function. The code attached to work can be used to compute the exact closed form of the expectations, altough the number of terms grows exponentially with the degree of the polynomial, making it impractical for high-degree polynomials.
    \item \emph{GeLU} \(\sigma(x) = x \Phi(x)\): this is a smooth approximation of the ReLU function, and it has been shown to have a closed form solution for the expectations \(I_2^\text{risk}\) and \(I_3\) \cite{richert2023layered}.
\end{itemize}
The other activations commonly used (such as sigmoid, tanh, etc.) do not have a closed-form solution for the expectations, and the dynamics can only be studied numerically by approximating the expectations.
One possibility is to use a Monte Carlo approximation of the expectations, but in my experience this is not very precise and often leads to worse results than simulating the SGD on weights directly. A better possibility is to approximate the activation function with Hermite polynomials, which have a closed-form solution for the expectations;
\addnote{Expanding in Hermite polynomials along every dimension is essentially equivalent to a tensor-product Gauss--Hermite quadrature of the correlated expectation; this technique is used in Chapter~\ref{chap:paper-2506}, in slightly different setting.}
a more advanced technique for approximating the expectations is discussed in \citet{citton2024line}, where the authors expand both the activation function and the correlated Gaussian measure in the Hermite basis to get a very precise approximation of the expectations for any activation function.

Despite the fact that the Saad \& Solla equations are not always analytically solvable, they are still a very powerful tool to study the dynamics of SGD around initialization. for small displacement, any \(sigma\) and \(h^\star\) can be approximated in a series of Hermite polynomials, and still have a closed form solution for the expectations, that approximate the dynamics of the order parameters. This is the core idea of many Chapters of this thesis, thus we will dig further in the next section.


\paragraph{Is \(d\to\infty\) all you need?}
When deriving Equations~\eqref{eq:intro:ss-odes}, we used as infinitesimal time step \(\sfrac{\gamma}{pd}\), which is indeed going to zero as \(d\to\infty\). However, in is not the only possibility to have a vanishing time step
\begin{itemize}
    \item \(\gamma \to 0\), \(d,p\) finite: this is the classical limit that for online SGD converges to \emph{gradient flow} \cite{robbins1951stochastic};
    \item \(p \to \infty\), \(d, \gamma\) finite: this is the limit of inifite width networks, sometimes called \emph{mean-field} limit, which has been studied in the context of two-layer networks \cite{mei2018mean, chizat2018global,sirignano2020mean}.
\end{itemize}
In Chapter~\ref{chap:paper-2302} we investigate the connection between the Saad \& Solla and this other two limits, covering them all with one single framework. Moreover, we study also possible combinations of those such as the \emph{high-dimensional gradient flow} limit, or the \emph{high-dimensional mean-field} limit, which are not covered by the classical Saad \& Solla equations.
% maybe expand more on the reslt of the paper 2302

